%%%%%%%%%%%%%%%%%%%%%%%%%%%%%%%%%%%%%%%%%%%%%%%%%%%%%%%%%%%%
% 7.1 ELEMENTARY PROBABILITY
% The Composition of Advanced Mathematics
%%%%%%%%%%%%%%%%%%%%%%%%%%%%%%%%%%%%%%%%%%%%%%%%%%%%%%%%%%%%

\section{Elementary Probability}
\label{sec:elementary-probability}

\prerequisite{
Sets, functions, basic counting, combinations, permutations, finite
sample spaces, and elementary algebra.
}

\learningobjectives{
By the end of this section, the reader should be able to:
\begin{itemize}
    \item explain probability as a mathematical model of uncertainty;
    \item define experiments, outcomes, sample spaces, and events;
    \item perform set operations on events;
    \item use the probability axioms;
    \item compute probabilities in finite equally likely sample spaces;
    \item use complements, unions, and intersections;
    \item apply the addition rule and inclusion--exclusion;
    \item distinguish mutually exclusive and independent events;
    \item compute conditional probabilities;
    \item apply the multiplication rule;
    \item apply the law of total probability;
    \item apply Bayes' theorem;
    \item use counting methods to solve probability problems;
    \item understand sampling with and without replacement;
    \item distinguish ordered and unordered experiments;
    \item use tree diagrams and probability tables;
    \item recognize odds and probability relationships;
    \item understand repeated independent trials;
    \item distinguish pairwise and mutual independence;
    \item understand the frequentist and classical interpretations of
          elementary probability;
    \item recognize the transition from finite probability to the more
          general axiomatic framework developed later in this chapter.
\end{itemize}
}

%%%%%%%%%%%%%%%%%%%%%%%%%%%%%%%%%%%%%%%%%%%%%%%%%%%%%%%%%%%%
\subsection{What Is Probability?}
\label{subsec:what-is-probability}
%%%%%%%%%%%%%%%%%%%%%%%%%%%%%%%%%%%%%%%%%%%%%%%%%%%%%%%%%%%%

Probability is the mathematics of uncertainty.

A deterministic mathematical model predicts one outcome from given
conditions.

A probabilistic model instead assigns numerical weights to possible
outcomes.

These weights describe how likely different events are to occur.

The central structure is

\[
\boxed{
\text{experiment}
\longrightarrow
\text{sample space}
\longrightarrow
\text{events}
\longrightarrow
\text{probabilities}.
}
\]

%%%%%%%%%%%%%%%%%%%%%%%%%%%%%%%%%%%%%%%%%%%%%%%%%%%%%%%%%%%%
\subsection{Random Experiments}
\label{subsec:random-experiments}
%%%%%%%%%%%%%%%%%%%%%%%%%%%%%%%%%%%%%%%%%%%%%%%%%%%%%%%%%%%%

A random experiment is a process whose outcome is not known in advance.

Examples include:

\begin{itemize}
    \item rolling a die;
    \item drawing a card;
    \item measuring whether a component fails;
    \item observing tomorrow's rainfall;
    \item recording the number of customers arriving in an hour.
\end{itemize}

The possible results of the experiment form the sample space.

%%%%%%%%%%%%%%%%%%%%%%%%%%%%%%%%%%%%%%%%%%%%%%%%%%%%%%%%%%%%
\subsection{Outcomes}
\label{subsec:probability-outcomes}
%%%%%%%%%%%%%%%%%%%%%%%%%%%%%%%%%%%%%%%%%%%%%%%%%%%%%%%%%%%%

An outcome is one possible result of an experiment.

For a six-sided die, possible outcomes are

\[
1,2,3,4,5,6.
\]

For two coin tosses, possible outcomes include

\[
HH,
\qquad
HT,
\qquad
TH,
\qquad
TT.
\]

An outcome is often denoted by the Greek letter

\[
\omega.
\]

The symbol

\[
\omega
\]

is called ``omega.''

%%%%%%%%%%%%%%%%%%%%%%%%%%%%%%%%%%%%%%%%%%%%%%%%%%%%%%%%%%%%
\subsection{Sample Spaces}
\label{subsec:sample-spaces}
%%%%%%%%%%%%%%%%%%%%%%%%%%%%%%%%%%%%%%%%%%%%%%%%%%%%%%%%%%%%

\begin{definitionbox}{Sample Space}
The sample space, usually denoted

\[
\boxed{
\Omega
}
\]

or \(S\), is the set of all possible outcomes of a random experiment.
\end{definitionbox}

The capital Greek letter

\[
\Omega
\]

is called ``capital omega.''

For one die roll,

\[
\boxed{
\Omega
=
\{1,2,3,4,5,6\}.
}
\]

%%%%%%%%%%%%%%%%%%%%%%%%%%%%%%%%%%%%%%%%%%%%%%%%%%%%%%%%%%%%
\subsection{Finite and Infinite Sample Spaces}
\label{subsec:finite-infinite-sample-spaces}
%%%%%%%%%%%%%%%%%%%%%%%%%%%%%%%%%%%%%%%%%%%%%%%%%%%%%%%%%%%%

A sample space may be finite.

For example,

\[
\Omega
=
\{H,T\}.
\]

It may be countably infinite.

For example, the number of coin tosses needed before the first head may
have sample space

\[
\boxed{
\Omega
=
\{1,2,3,\ldots\}.
}
\]

It may also be uncountable.

For example, a measured lifetime may take any value in

\[
[0,\infty).
\]

The mathematical treatment becomes increasingly sophisticated as the
sample space becomes more general.

%%%%%%%%%%%%%%%%%%%%%%%%%%%%%%%%%%%%%%%%%%%%%%%%%%%%%%%%%%%%
\subsection{Events}
\label{subsec:probability-events}
%%%%%%%%%%%%%%%%%%%%%%%%%%%%%%%%%%%%%%%%%%%%%%%%%%%%%%%%%%%%

\begin{definitionbox}{Event}
An event is a subset of the sample space.
\end{definitionbox}

Thus,

\[
\boxed{
A\subseteq\Omega.
}
\]

An event occurs when the observed outcome belongs to \(A\).

For a die roll, the event

\[
A=\{2,4,6\}
\]

means ``the result is even.''

%%%%%%%%%%%%%%%%%%%%%%%%%%%%%%%%%%%%%%%%%%%%%%%%%%%%%%%%%%%%
\subsection{Simple and Compound Events}
\label{subsec:simple-compound-events}
%%%%%%%%%%%%%%%%%%%%%%%%%%%%%%%%%%%%%%%%%%%%%%%%%%%%%%%%%%%%

A simple event contains one outcome.

For example,

\[
A=\{4\}.
\]

A compound event contains several outcomes.

For example,

\[
B=\{2,4,6\}.
\]

The impossible event is

\[
\boxed{
\varnothing.
}
\]

The certain event is

\[
\boxed{
\Omega.
}
\]

%%%%%%%%%%%%%%%%%%%%%%%%%%%%%%%%%%%%%%%%%%%%%%%%%%%%%%%%%%%%
\subsection{Set Operations on Events}
\label{subsec:set-operations-events}
%%%%%%%%%%%%%%%%%%%%%%%%%%%%%%%%%%%%%%%%%%%%%%%%%%%%%%%%%%%%

Because events are sets, ordinary set operations become probability
operations.

For events \(A\) and \(B\):

\[
A\cup B
\]

means ``\(A\) or \(B\) occurs.''

\[
A\cap B
\]

means ``both \(A\) and \(B\) occur.''

\[
A^c
\]

means ``\(A\) does not occur.''

\[
A\setminus B
\]

means ``\(A\) occurs but \(B\) does not.''

%%%%%%%%%%%%%%%%%%%%%%%%%%%%%%%%%%%%%%%%%%%%%%%%%%%%%%%%%%%%
\subsection{Union of Events}
\label{subsec:event-union}
%%%%%%%%%%%%%%%%%%%%%%%%%%%%%%%%%%%%%%%%%%%%%%%%%%%%%%%%%%%%

The union

\[
\boxed{
A\cup B
}
\]

contains outcomes belonging to \(A\), \(B\), or both.

In probability, the word ``or'' is usually inclusive.

Thus,

\[
A\cup B
\]

means at least one of the events occurs.

%%%%%%%%%%%%%%%%%%%%%%%%%%%%%%%%%%%%%%%%%%%%%%%%%%%%%%%%%%%%
\subsection{Intersection of Events}
\label{subsec:event-intersection}
%%%%%%%%%%%%%%%%%%%%%%%%%%%%%%%%%%%%%%%%%%%%%%%%%%%%%%%%%%%%

The intersection

\[
\boxed{
A\cap B
}
\]

contains outcomes that belong to both events.

Thus,

\[
A\cap B
\]

represents simultaneous occurrence.

%%%%%%%%%%%%%%%%%%%%%%%%%%%%%%%%%%%%%%%%%%%%%%%%%%%%%%%%%%%%
\subsection{Complement of an Event}
\label{subsec:event-complement}
%%%%%%%%%%%%%%%%%%%%%%%%%%%%%%%%%%%%%%%%%%%%%%%%%%%%%%%%%%%%

The complement

\[
\boxed{
A^c
}
\]

contains every outcome in \(\Omega\) that is not in \(A\).

Therefore,

\[
A\cup A^c=\Omega
\]

and

\[
A\cap A^c=\varnothing.
\]

%%%%%%%%%%%%%%%%%%%%%%%%%%%%%%%%%%%%%%%%%%%%%%%%%%%%%%%%%%%%
\subsection{De Morgan's Laws for Events}
\label{subsec:demorgan-probability}
%%%%%%%%%%%%%%%%%%%%%%%%%%%%%%%%%%%%%%%%%%%%%%%%%%%%%%%%%%%%

De Morgan's laws give

\[
\boxed{
(A\cup B)^c
=
A^c\cap B^c
}
\]

and

\[
\boxed{
(A\cap B)^c
=
A^c\cup B^c.
}
\]

Thus

\[
\text{``neither \(A\) nor \(B\)''}
\]

corresponds to

\[
A^c\cap B^c.
\]

%%%%%%%%%%%%%%%%%%%%%%%%%%%%%%%%%%%%%%%%%%%%%%%%%%%%%%%%%%%%
\subsection{Probability Function}
\label{subsec:probability-function}
%%%%%%%%%%%%%%%%%%%%%%%%%%%%%%%%%%%%%%%%%%%%%%%%%%%%%%%%%%%%

A probability function assigns a number to each event.

We write

\[
\boxed{
P(A)
}
\]

for the probability of event \(A\).

The notation is read

\[
\text{``the probability of \(A\).''}
\]

Probabilities satisfy fundamental axioms.

%%%%%%%%%%%%%%%%%%%%%%%%%%%%%%%%%%%%%%%%%%%%%%%%%%%%%%%%%%%%
\subsection{Kolmogorov Probability Axioms}
\label{subsec:probability-axioms}
%%%%%%%%%%%%%%%%%%%%%%%%%%%%%%%%%%%%%%%%%%%%%%%%%%%%%%%%%%%%

\begin{definitionbox}{Probability Measure}
A probability measure \(P\) satisfies:

\begin{enumerate}
    \item Nonnegativity:
    \[
    \boxed{
    P(A)\geq0
    }
    \]

    for every event \(A\).

    \item Normalization:
    \[
    \boxed{
    P(\Omega)=1.
    }
    \]

    \item Countable additivity:
    if
    \[
    A_1,A_2,\ldots
    \]
    are pairwise disjoint, then
    \[
    \boxed{
    P\left(
    \bigcup_{i=1}^{\infty}A_i
    \right)
    =
    \sum_{i=1}^{\infty}P(A_i).
    }
    \]
\end{enumerate}
\end{definitionbox}

These are the Kolmogorov axioms of probability.

%%%%%%%%%%%%%%%%%%%%%%%%%%%%%%%%%%%%%%%%%%%%%%%%%%%%%%%%%%%%
\subsection{Probability Is Between Zero and One}
\label{subsec:probability-range}
%%%%%%%%%%%%%%%%%%%%%%%%%%%%%%%%%%%%%%%%%%%%%%%%%%%%%%%%%%%%

For every event \(A\),

\[
\boxed{
0\leq P(A)\leq1.
}
\]

The lower bound follows from nonnegativity.

For the upper bound, write

\[
\Omega=A\cup A^c,
\]

where the union is disjoint.

Then

\[
1
=
P(\Omega)
=
P(A)+P(A^c).
\]

Since

\[
P(A^c)\geq0,
\]

we obtain

\[
P(A)\leq1.
\]

%%%%%%%%%%%%%%%%%%%%%%%%%%%%%%%%%%%%%%%%%%%%%%%%%%%%%%%%%%%%
\subsection{Probability of the Empty Event}
\label{subsec:empty-event-probability}
%%%%%%%%%%%%%%%%%%%%%%%%%%%%%%%%%%%%%%%%%%%%%%%%%%%%%%%%%%%%

The impossible event has probability

\[
\boxed{
P(\varnothing)=0.
}
\]

To see this,

\[
\Omega
=
\Omega\cup\varnothing
\]

as a disjoint union.

Therefore,

\[
P(\Omega)
=
P(\Omega)+P(\varnothing),
\]

so

\[
P(\varnothing)=0.
\]

%%%%%%%%%%%%%%%%%%%%%%%%%%%%%%%%%%%%%%%%%%%%%%%%%%%%%%%%%%%%
\subsection{Complement Rule}
\label{subsec:complement-rule}
%%%%%%%%%%%%%%%%%%%%%%%%%%%%%%%%%%%%%%%%%%%%%%%%%%%%%%%%%%%%

Since

\[
A\cup A^c=\Omega
\]

and the two events are disjoint,

\[
P(A)+P(A^c)=1.
\]

Therefore,

\[
\boxed{
P(A^c)=1-P(A).
}
\]

The complement rule is often the easiest way to calculate
``at least one'' probabilities.

%%%%%%%%%%%%%%%%%%%%%%%%%%%%%%%%%%%%%%%%%%%%%%%%%%%%%%%%%%%%
\subsection{Worked Example: Complement Rule}
\label{subsec:worked-complement-rule}
%%%%%%%%%%%%%%%%%%%%%%%%%%%%%%%%%%%%%%%%%%%%%%%%%%%%%%%%%%%%

\begin{workedexamplebox}{At Least One Six}

Roll a fair die four times.

Find the probability of obtaining at least one \(6\).

The complement event is

\[
\text{no six occurs}.
\]

For one roll,

\[
P(\text{not }6)=\frac56.
\]

Assuming independent rolls,

\[
P(\text{no six in four rolls})
=
\left(\frac56\right)^4.
\]

Therefore,

\begin{answerbox}
\[
\boxed{
P(\text{at least one }6)
=
1-\left(\frac56\right)^4.
}
\]
\end{answerbox}

\end{workedexamplebox}

%%%%%%%%%%%%%%%%%%%%%%%%%%%%%%%%%%%%%%%%%%%%%%%%%%%%%%%%%%%%
\subsection{Monotonicity}
\label{subsec:probability-monotonicity}
%%%%%%%%%%%%%%%%%%%%%%%%%%%%%%%%%%%%%%%%%%%%%%%%%%%%%%%%%%%%

If

\[
A\subseteq B,
\]

then

\[
\boxed{
P(A)\leq P(B).
}
\]

Indeed,

\[
B
=
A\cup(B\setminus A),
\]

where the union is disjoint.

Therefore,

\[
P(B)
=
P(A)+P(B\setminus A)
\geq
P(A).
\]

%%%%%%%%%%%%%%%%%%%%%%%%%%%%%%%%%%%%%%%%%%%%%%%%%%%%%%%%%%%%
\subsection{Finite Equally Likely Outcomes}
\label{subsec:equally-likely-outcomes}
%%%%%%%%%%%%%%%%%%%%%%%%%%%%%%%%%%%%%%%%%%%%%%%%%%%%%%%%%%%%

Suppose the finite sample space contains

\[
N
\]

equally likely outcomes.

Then each outcome has probability

\[
\frac1N.
\]

If event \(A\) contains \(m\) outcomes, then

\[
\boxed{
P(A)
=
\frac{|A|}{|\Omega|}
=
\frac{m}{N}.
}
\]

This is the classical counting formula for probability.

%%%%%%%%%%%%%%%%%%%%%%%%%%%%%%%%%%%%%%%%%%%%%%%%%%%%%%%%%%%%
\subsection{Worked Example: Fair Die}
\label{subsec:worked-fair-die}
%%%%%%%%%%%%%%%%%%%%%%%%%%%%%%%%%%%%%%%%%%%%%%%%%%%%%%%%%%%%

\begin{workedexamplebox}{Rolling an Even Number}

Roll a fair six-sided die.

The sample space is

\[
\Omega
=
\{1,2,3,4,5,6\}.
\]

Let

\[
A=\{2,4,6\}.
\]

There are \(3\) favorable outcomes among \(6\) equally likely outcomes.

\begin{answerbox}
\[
\boxed{
P(A)
=
\frac36
=
\frac12.
}
\]
\end{answerbox}

\end{workedexamplebox}

%%%%%%%%%%%%%%%%%%%%%%%%%%%%%%%%%%%%%%%%%%%%%%%%%%%%%%%%%%%%
\subsection{Addition Rule}
\label{subsec:addition-rule-probability}
%%%%%%%%%%%%%%%%%%%%%%%%%%%%%%%%%%%%%%%%%%%%%%%%%%%%%%%%%%%%

For any two events \(A\) and \(B\),

\[
\boxed{
P(A\cup B)
=
P(A)+P(B)-P(A\cap B).
}
\]

The intersection term is subtracted because it is counted twice in

\[
P(A)+P(B).
\]

This is the two-event inclusion--exclusion formula.

%%%%%%%%%%%%%%%%%%%%%%%%%%%%%%%%%%%%%%%%%%%%%%%%%%%%%%%%%%%%
\subsection{Mutually Exclusive Events}
\label{subsec:mutually-exclusive-events}
%%%%%%%%%%%%%%%%%%%%%%%%%%%%%%%%%%%%%%%%%%%%%%%%%%%%%%%%%%%%

\begin{definitionbox}{Mutually Exclusive Events}
Events \(A\) and \(B\) are mutually exclusive, or disjoint, if

\[
\boxed{
A\cap B=\varnothing.
}
\]
\end{definitionbox}

They cannot occur simultaneously.

For disjoint events,

\[
\boxed{
P(A\cup B)
=
P(A)+P(B).
}
\]

%%%%%%%%%%%%%%%%%%%%%%%%%%%%%%%%%%%%%%%%%%%%%%%%%%%%%%%%%%%%
\subsection{Worked Example: Addition Rule}
\label{subsec:worked-addition-rule}
%%%%%%%%%%%%%%%%%%%%%%%%%%%%%%%%%%%%%%%%%%%%%%%%%%%%%%%%%%%%

\begin{workedexamplebox}{Card Draw}

Draw one card from a standard \(52\)-card deck.

Let

\[
A
=
\text{the card is a heart},
\]

and

\[
B
=
\text{the card is a king}.
\]

Then

\[
P(A)=\frac{13}{52},
\]

\[
P(B)=\frac4{52},
\]

and

\[
P(A\cap B)=\frac1{52}
\]

because the king of hearts belongs to both events.

Thus,

\[
P(A\cup B)
=
\frac{13}{52}
+
\frac4{52}
-
\frac1{52}.
\]

\begin{answerbox}
\[
\boxed{
P(A\cup B)
=
\frac{16}{52}
=
\frac4{13}.
}
\]
\end{answerbox}

\end{workedexamplebox}

%%%%%%%%%%%%%%%%%%%%%%%%%%%%%%%%%%%%%%%%%%%%%%%%%%%%%%%%%%%%
\subsection{Three-Event Inclusion--Exclusion}
\label{subsec:three-event-inclusion-exclusion}
%%%%%%%%%%%%%%%%%%%%%%%%%%%%%%%%%%%%%%%%%%%%%%%%%%%%%%%%%%%%

For events \(A,B,C\),

\[
\boxed{
\begin{aligned}
P(A\cup B\cup C)
={}&
P(A)+P(B)+P(C)
\\
&-
P(A\cap B)
-
P(A\cap C)
-
P(B\cap C)
\\
&+
P(A\cap B\cap C).
\end{aligned}
}
\]

This is the probability form of inclusion--exclusion.

%%%%%%%%%%%%%%%%%%%%%%%%%%%%%%%%%%%%%%%%%%%%%%%%%%%%%%%%%%%%
\subsection{Conditional Probability}
\label{subsec:conditional-probability}
%%%%%%%%%%%%%%%%%%%%%%%%%%%%%%%%%%%%%%%%%%%%%%%%%%%%%%%%%%%%

Suppose event \(B\) is known to have occurred.

The relevant sample space is then reduced from \(\Omega\) to \(B\).

\begin{definitionbox}{Conditional Probability}
If

\[
P(B)>0,
\]

the conditional probability of \(A\) given \(B\) is

\[
\boxed{
P(A\mid B)
=
\frac{P(A\cap B)}{P(B)}.
}
\]
\end{definitionbox}

The notation

\[
P(A\mid B)
\]

is read

\[
\text{``the probability of \(A\) given \(B\).''}
\]

%%%%%%%%%%%%%%%%%%%%%%%%%%%%%%%%%%%%%%%%%%%%%%%%%%%%%%%%%%%%
\subsection{Worked Example: Conditional Card Probability}
\label{subsec:worked-conditional-card}
%%%%%%%%%%%%%%%%%%%%%%%%%%%%%%%%%%%%%%%%%%%%%%%%%%%%%%%%%%%%

\begin{workedexamplebox}{Given That the Card Is a Face Card}

Draw one card from a standard deck.

Let

\[
A
=
\text{the card is a king},
\]

and

\[
B
=
\text{the card is a face card}.
\]

There are \(12\) face cards:

\[
4\text{ jacks},
\qquad
4\text{ queens},
\qquad
4\text{ kings}.
\]

Thus,

\[
P(B)=\frac{12}{52}.
\]

Every king is a face card, so

\[
P(A\cap B)=\frac4{52}.
\]

Therefore,

\[
P(A\mid B)
=
\frac{4/52}{12/52}.
\]

\begin{answerbox}
\[
\boxed{
P(A\mid B)=\frac13.
}
\]
\end{answerbox}

\end{workedexamplebox}

%%%%%%%%%%%%%%%%%%%%%%%%%%%%%%%%%%%%%%%%%%%%%%%%%%%%%%%%%%%%
\subsection{Multiplication Rule}
\label{subsec:multiplication-rule}
%%%%%%%%%%%%%%%%%%%%%%%%%%%%%%%%%%%%%%%%%%%%%%%%%%%%%%%%%%%%

Rearranging the conditional probability formula gives

\[
\boxed{
P(A\cap B)
=
P(A\mid B)P(B).
}
\]

Similarly,

\[
\boxed{
P(A\cap B)
=
P(B\mid A)P(A).
}
\]

These are multiplication rules for intersections.

%%%%%%%%%%%%%%%%%%%%%%%%%%%%%%%%%%%%%%%%%%%%%%%%%%%%%%%%%%%%
\subsection{Sequential Multiplication Rule}
\label{subsec:sequential-multiplication-rule}
%%%%%%%%%%%%%%%%%%%%%%%%%%%%%%%%%%%%%%%%%%%%%%%%%%%%%%%%%%%%

For events

\[
A_1,\ldots,A_n,
\]

we have

\[
\boxed{
\begin{aligned}
P(A_1\cap\cdots\cap A_n)
={}&
P(A_1)
P(A_2\mid A_1)
\\
&\cdot
P(A_3\mid A_1\cap A_2)
\cdots
\\
&\cdot
P(A_n\mid A_1\cap\cdots\cap A_{n-1}).
\end{aligned}
}
\]

This is sometimes called the chain rule for probability.

%%%%%%%%%%%%%%%%%%%%%%%%%%%%%%%%%%%%%%%%%%%%%%%%%%%%%%%%%%%%
\subsection{Sampling Without Replacement}
\label{subsec:sampling-without-replacement}
%%%%%%%%%%%%%%%%%%%%%%%%%%%%%%%%%%%%%%%%%%%%%%%%%%%%%%%%%%%%

Suppose an urn contains

\[
5
\]

red balls and

\[
3
\]

blue balls.

Two balls are selected without replacement.

The probability the first is red is

\[
\frac58.
\]

After a red ball is removed, the probability the second is red becomes

\[
\frac47.
\]

Therefore,

\[
\boxed{
P(\text{two red})
=
\frac58\cdot\frac47
=
\frac5{14}.
}
\]

The second probability changes because the first selection alters the
composition of the urn.

%%%%%%%%%%%%%%%%%%%%%%%%%%%%%%%%%%%%%%%%%%%%%%%%%%%%%%%%%%%%
\subsection{Sampling With Replacement}
\label{subsec:sampling-with-replacement}
%%%%%%%%%%%%%%%%%%%%%%%%%%%%%%%%%%%%%%%%%%%%%%%%%%%%%%%%%%%%

If the first selected object is returned before the second selection,
the original probabilities are restored.

Using the same urn,

\[
P(\text{red on first})
=
\frac58,
\]

and

\[
P(\text{red on second})
=
\frac58.
\]

Hence,

\[
\boxed{
P(\text{two red with replacement})
=
\left(\frac58\right)^2.
}
\]

Replacement often creates independence between repeated selections.

%%%%%%%%%%%%%%%%%%%%%%%%%%%%%%%%%%%%%%%%%%%%%%%%%%%%%%%%%%%%
\subsection{Independent Events}
\label{subsec:independent-events}
%%%%%%%%%%%%%%%%%%%%%%%%%%%%%%%%%%%%%%%%%%%%%%%%%%%%%%%%%%%%

\begin{definitionbox}{Independent Events}
Events \(A\) and \(B\) are independent if

\[
\boxed{
P(A\cap B)
=
P(A)P(B).
}
\]
\end{definitionbox}

If both probabilities are positive, this is equivalent to

\[
\boxed{
P(A\mid B)=P(A)
}
\]

and

\[
\boxed{
P(B\mid A)=P(B).
}
\]

Thus knowing one event occurred does not change the probability of the
other.

%%%%%%%%%%%%%%%%%%%%%%%%%%%%%%%%%%%%%%%%%%%%%%%%%%%%%%%%%%%%
\subsection{Worked Example: Independent Coin Tosses}
\label{subsec:worked-independent-coins}
%%%%%%%%%%%%%%%%%%%%%%%%%%%%%%%%%%%%%%%%%%%%%%%%%%%%%%%%%%%%

\begin{workedexamplebox}{Two Fair Tosses}

Toss a fair coin twice.

Let

\[
A
=
\text{first toss is heads}
\]

and

\[
B
=
\text{second toss is heads}.
\]

Then

\[
P(A)=\frac12,
\qquad
P(B)=\frac12.
\]

Also,

\[
P(A\cap B)=\frac14.
\]

Since

\[
\frac14
=
\frac12\cdot\frac12,
\]

\begin{answerbox}
\[
\boxed{
A\text{ and }B\text{ are independent}.
}
\]
\end{answerbox}

\end{workedexamplebox}

%%%%%%%%%%%%%%%%%%%%%%%%%%%%%%%%%%%%%%%%%%%%%%%%%%%%%%%%%%%%
\subsection{Mutually Exclusive Versus Independent}
\label{subsec:exclusive-versus-independent}
%%%%%%%%%%%%%%%%%%%%%%%%%%%%%%%%%%%%%%%%%%%%%%%%%%%%%%%%%%%%

These concepts are very different.

Mutually exclusive means

\[
\boxed{
P(A\cap B)=0.
}
\]

Independent means

\[
\boxed{
P(A\cap B)=P(A)P(B).
}
\]

If \(A\) and \(B\) are mutually exclusive and both have positive
probability, then

\[
P(A)P(B)>0,
\]

so they cannot be independent.

\begin{warningbox}
Mutually exclusive events with positive probability are dependent, not
independent.
\end{warningbox}

%%%%%%%%%%%%%%%%%%%%%%%%%%%%%%%%%%%%%%%%%%%%%%%%%%%%%%%%%%%%
\subsection{Pairwise Independence}
\label{subsec:pairwise-independence}
%%%%%%%%%%%%%%%%%%%%%%%%%%%%%%%%%%%%%%%%%%%%%%%%%%%%%%%%%%%%

Events

\[
A_1,\ldots,A_n
\]

are pairwise independent if every distinct pair satisfies

\[
\boxed{
P(A_i\cap A_j)
=
P(A_i)P(A_j).
}
\]

Pairwise independence does not necessarily imply mutual independence.

%%%%%%%%%%%%%%%%%%%%%%%%%%%%%%%%%%%%%%%%%%%%%%%%%%%%%%%%%%%%
\subsection{Mutual Independence}
\label{subsec:mutual-independence}
%%%%%%%%%%%%%%%%%%%%%%%%%%%%%%%%%%%%%%%%%%%%%%%%%%%%%%%%%%%%

Events

\[
A_1,\ldots,A_n
\]

are mutually independent if for every nonempty index set

\[
J\subseteq\{1,\ldots,n\},
\]

we have

\[
\boxed{
P\left(
\bigcap_{j\in J}A_j
\right)
=
\prod_{j\in J}P(A_j).
}
\]

Thus every possible finite subcollection must factor.

%%%%%%%%%%%%%%%%%%%%%%%%%%%%%%%%%%%%%%%%%%%%%%%%%%%%%%%%%%%%
\subsection{Pairwise but Not Mutually Independent Example}
\label{subsec:pairwise-not-mutual-example}
%%%%%%%%%%%%%%%%%%%%%%%%%%%%%%%%%%%%%%%%%%%%%%%%%%%%%%%%%%%%

Toss two fair coins.

Let

\[
A
=
\text{first coin is heads},
\]

\[
B
=
\text{second coin is heads},
\]

and

\[
C
=
\text{the two coins show the same result}.
\]

Each event has probability

\[
\frac12.
\]

Also,

\[
P(A\cap B)
=
P(A\cap C)
=
P(B\cap C)
=
\frac14.
\]

Thus each pair is independent.

However,

\[
A\cap B\cap C
=
\{HH\},
\]

so

\[
P(A\cap B\cap C)=\frac14.
\]

But

\[
P(A)P(B)P(C)
=
\frac18.
\]

Therefore the three events are not mutually independent.

%%%%%%%%%%%%%%%%%%%%%%%%%%%%%%%%%%%%%%%%%%%%%%%%%%%%%%%%%%%%
\subsection{Partitions of the Sample Space}
\label{subsec:sample-space-partitions}
%%%%%%%%%%%%%%%%%%%%%%%%%%%%%%%%%%%%%%%%%%%%%%%%%%%%%%%%%%%%

A collection of events

\[
B_1,\ldots,B_n
\]

forms a partition of \(\Omega\) if:

\[
B_i\cap B_j=\varnothing
\]

for

\[
i\neq j,
\]

and

\[
\boxed{
\bigcup_{i=1}^{n}B_i=\Omega.
}
\]

Every outcome belongs to exactly one part.

Partitions are essential for the law of total probability.

%%%%%%%%%%%%%%%%%%%%%%%%%%%%%%%%%%%%%%%%%%%%%%%%%%%%%%%%%%%%
\subsection{Law of Total Probability}
\label{subsec:law-total-probability}
%%%%%%%%%%%%%%%%%%%%%%%%%%%%%%%%%%%%%%%%%%%%%%%%%%%%%%%%%%%%

\begin{theorem}[Law of Total Probability]
Suppose

\[
B_1,\ldots,B_n
\]

form a partition of \(\Omega\) with

\[
P(B_i)>0.
\]

Then for every event \(A\),

\[
\boxed{
P(A)
=
\sum_{i=1}^{n}
P(A\mid B_i)P(B_i).
}
\]
\end{theorem}

The event \(A\) is decomposed according to the mutually exclusive cases
\(B_i\).

%%%%%%%%%%%%%%%%%%%%%%%%%%%%%%%%%%%%%%%%%%%%%%%%%%%%%%%%%%%%
\subsection{Worked Example: Total Probability}
\label{subsec:worked-total-probability}
%%%%%%%%%%%%%%%%%%%%%%%%%%%%%%%%%%%%%%%%%%%%%%%%%%%%%%%%%%%%

\begin{workedexamplebox}{Two Production Lines}

Suppose a factory has two production lines.

Line \(1\) produces

\[
60\%
\]

of all items and has defect probability

\[
0.02.
\]

Line \(2\) produces

\[
40\%
\]

of all items and has defect probability

\[
0.05.
\]

Let \(D\) be the event that an item is defective.

Then

\[
P(D)
=
P(D\mid L_1)P(L_1)
+
P(D\mid L_2)P(L_2).
\]

Thus,

\[
P(D)
=
(0.02)(0.60)
+
(0.05)(0.40).
\]

\begin{answerbox}
\[
\boxed{
P(D)=0.032.
}
\]
\end{answerbox}

\end{workedexamplebox}

%%%%%%%%%%%%%%%%%%%%%%%%%%%%%%%%%%%%%%%%%%%%%%%%%%%%%%%%%%%%
\subsection{Bayes' Theorem}
\label{subsec:bayes-theorem}
%%%%%%%%%%%%%%%%%%%%%%%%%%%%%%%%%%%%%%%%%%%%%%%%%%%%%%%%%%%%

Bayes' theorem reverses a conditional probability.

\begin{theorem}[Bayes' Theorem]
For events \(A\) and \(B\) with positive probabilities,

\[
\boxed{
P(B\mid A)
=
\frac{
P(A\mid B)P(B)
}{
P(A)
}.
}
\]
\end{theorem}

If

\[
B_1,\ldots,B_n
\]

form a partition, then

\[
\boxed{
P(B_j\mid A)
=
\frac{
P(A\mid B_j)P(B_j)
}{
\sum_{i=1}^{n}
P(A\mid B_i)P(B_i)
}.
}
\]

%%%%%%%%%%%%%%%%%%%%%%%%%%%%%%%%%%%%%%%%%%%%%%%%%%%%%%%%%%%%
\subsection{Prior and Posterior Probabilities}
\label{subsec:prior-posterior}
%%%%%%%%%%%%%%%%%%%%%%%%%%%%%%%%%%%%%%%%%%%%%%%%%%%%%%%%%%%%

In Bayesian terminology,

\[
P(B)
\]

is the prior probability.

The quantity

\[
P(A\mid B)
\]

is the likelihood of observing \(A\) under \(B\).

The quantity

\[
P(B\mid A)
\]

is the posterior probability after observing \(A\).

Thus Bayes' theorem updates uncertainty using evidence.

%%%%%%%%%%%%%%%%%%%%%%%%%%%%%%%%%%%%%%%%%%%%%%%%%%%%%%%%%%%%
\subsection{Worked Example: Bayes' Theorem}
\label{subsec:worked-bayes-production}
%%%%%%%%%%%%%%%%%%%%%%%%%%%%%%%%%%%%%%%%%%%%%%%%%%%%%%%%%%%%

\begin{workedexamplebox}{Which Production Line Made the Defective Item?}

Continue the production-line example.

We know

\[
P(L_1)=0.60,
\]

\[
P(D\mid L_1)=0.02,
\]

and

\[
P(D)=0.032.
\]

Then

\[
P(L_1\mid D)
=
\frac{
P(D\mid L_1)P(L_1)
}{
P(D)
}.
\]

Therefore,

\[
P(L_1\mid D)
=
\frac{(0.02)(0.60)}{0.032}.
\]

\begin{answerbox}
\[
\boxed{
P(L_1\mid D)=0.375.
}
\]
\end{answerbox}

Thus \(37.5\%\) of defective items are expected to have come from
line \(1\).

\end{workedexamplebox}

%%%%%%%%%%%%%%%%%%%%%%%%%%%%%%%%%%%%%%%%%%%%%%%%%%%%%%%%%%%%
\subsection{Bayes' Theorem and Diagnostic Testing}
\label{subsec:bayes-diagnostic-testing}
%%%%%%%%%%%%%%%%%%%%%%%%%%%%%%%%%%%%%%%%%%%%%%%%%%%%%%%%%%%%

Bayes' theorem is especially important when interpreting tests.

Suppose:

\[
D
=
\text{condition is present},
\]

and

\[
+
=
\text{test is positive}.
\]

Then

\[
P(+\mid D)
\]

is not the same as

\[
P(D\mid +).
\]

The first is the probability of a positive test given the condition.

The second is the probability the condition is present given a positive
test.

Confusing these quantities is a common error.

%%%%%%%%%%%%%%%%%%%%%%%%%%%%%%%%%%%%%%%%%%%%%%%%%%%%%%%%%%%%
\subsection{Sensitivity and Specificity}
\label{subsec:sensitivity-specificity}
%%%%%%%%%%%%%%%%%%%%%%%%%%%%%%%%%%%%%%%%%%%%%%%%%%%%%%%%%%%%

For a binary diagnostic test:

\[
\boxed{
\text{sensitivity}
=
P(+\mid D)
}
\]

and

\[
\boxed{
\text{specificity}
=
P(-\mid D^c).
}
\]

The false-positive probability is

\[
\boxed{
P(+\mid D^c)
=
1-\text{specificity}.
}
\]

The predictive probability

\[
P(D\mid+)
\]

also depends strongly on the prevalence

\[
P(D).
\]

%%%%%%%%%%%%%%%%%%%%%%%%%%%%%%%%%%%%%%%%%%%%%%%%%%%%%%%%%%%%
\subsection{Worked Example: Rare Condition}
\label{subsec:worked-rare-condition}
%%%%%%%%%%%%%%%%%%%%%%%%%%%%%%%%%%%%%%%%%%%%%%%%%%%%%%%%%%%%

\begin{workedexamplebox}{Positive Test for a Rare Condition}

Suppose

\[
P(D)=0.01,
\]

\[
P(+\mid D)=0.95,
\]

and

\[
P(+\mid D^c)=0.05.
\]

By total probability,

\[
P(+)
=
(0.95)(0.01)
+
(0.05)(0.99).
\]

Thus,

\[
P(+)
=
0.0095+0.0495
=
0.059.
\]

Bayes' theorem gives

\[
P(D\mid+)
=
\frac{(0.95)(0.01)}{0.059}.
\]

\begin{answerbox}
\[
\boxed{
P(D\mid+)
\approx0.161.
}
\]
\end{answerbox}

Despite high sensitivity, the posterior probability is only about
\(16.1\%\) because the condition is rare.

\end{workedexamplebox}

%%%%%%%%%%%%%%%%%%%%%%%%%%%%%%%%%%%%%%%%%%%%%%%%%%%%%%%%%%%%
\subsection{Counting and Probability}
\label{subsec:counting-and-probability}
%%%%%%%%%%%%%%%%%%%%%%%%%%%%%%%%%%%%%%%%%%%%%%%%%%%%%%%%%%%%

In finite equally likely experiments,

\[
P(A)
=
\frac{|A|}{|\Omega|}.
\]

Therefore probability calculations often reduce to combinatorial
counting.

Important tools include:

\[
\boxed{
\text{permutations},
\quad
\text{combinations},
\quad
\text{multinomial coefficients}.
}
\]

%%%%%%%%%%%%%%%%%%%%%%%%%%%%%%%%%%%%%%%%%%%%%%%%%%%%%%%%%%%%
\subsection{Ordered Sampling}
\label{subsec:ordered-sampling-probability}
%%%%%%%%%%%%%%%%%%%%%%%%%%%%%%%%%%%%%%%%%%%%%%%%%%%%%%%%%%%%

Suppose \(k\) distinct objects are selected in order from \(n\) distinct
objects without replacement.

The number of possible ordered outcomes is

\[
\boxed{
n(n-1)\cdots(n-k+1)
=
\frac{n!}{(n-k)!}.
}
\]

This is a permutation count.

%%%%%%%%%%%%%%%%%%%%%%%%%%%%%%%%%%%%%%%%%%%%%%%%%%%%%%%%%%%%
\subsection{Unordered Sampling}
\label{subsec:unordered-sampling-probability}
%%%%%%%%%%%%%%%%%%%%%%%%%%%%%%%%%%%%%%%%%%%%%%%%%%%%%%%%%%%%

If order does not matter, the number of \(k\)-element samples is

\[
\boxed{
\binom nk.
}
\]

A common probability error occurs when the numerator and denominator
use incompatible conventions about order.

%%%%%%%%%%%%%%%%%%%%%%%%%%%%%%%%%%%%%%%%%%%%%%%%%%%%%%%%%%%%
\subsection{Worked Example: Poker Hand}
\label{subsec:worked-poker-hand}
%%%%%%%%%%%%%%%%%%%%%%%%%%%%%%%%%%%%%%%%%%%%%%%%%%%%%%%%%%%%

\begin{workedexamplebox}{Probability of Four Aces in a Five-Card Hand}

A five-card poker hand is an unordered selection.

The total number of hands is

\[
\binom{52}{5}.
\]

To have all four aces, choose:

\[
4
\]

aces from the \(4\) available and

\[
1
\]

additional card from the \(48\) non-aces.

Thus the favorable count is

\[
\binom44\binom{48}{1}.
\]

Therefore,

\begin{answerbox}
\[
\boxed{
P(\text{four aces})
=
\frac{
\binom44\binom{48}{1}
}{
\binom{52}{5}
}.
}
\]
\end{answerbox}

\end{workedexamplebox}

%%%%%%%%%%%%%%%%%%%%%%%%%%%%%%%%%%%%%%%%%%%%%%%%%%%%%%%%%%%%
\subsection{Hypergeometric Probability}
\label{subsec:hypergeometric-probability-preview}
%%%%%%%%%%%%%%%%%%%%%%%%%%%%%%%%%%%%%%%%%%%%%%%%%%%%%%%%%%%%

Suppose a population contains:

\[
N
\]

objects,

\[
K
\]

of which are classified as successes.

Choose

\[
n
\]

objects without replacement.

The probability of obtaining exactly \(k\) successes is

\[
\boxed{
P(X=k)
=
\frac{
\binom{K}{k}
\binom{N-K}{n-k}
}{
\binom{N}{n}
}.
}
\]

This becomes the hypergeometric distribution, developed later in the
chapter.

%%%%%%%%%%%%%%%%%%%%%%%%%%%%%%%%%%%%%%%%%%%%%%%%%%%%%%%%%%%%
\subsection{Repeated Independent Trials}
\label{subsec:repeated-independent-trials}
%%%%%%%%%%%%%%%%%%%%%%%%%%%%%%%%%%%%%%%%%%%%%%%%%%%%%%%%%%%%

Suppose an experiment is repeated independently.

If success occurs with probability

\[
p,
\]

then failure occurs with probability

\[
1-p.
\]

A specified sequence containing \(k\) successes and \(n-k\) failures
has probability

\[
\boxed{
p^k(1-p)^{n-k}.
}
\]

%%%%%%%%%%%%%%%%%%%%%%%%%%%%%%%%%%%%%%%%%%%%%%%%%%%%%%%%%%%%
\subsection{Exactly \(k\) Successes}
\label{subsec:exactly-k-successes}
%%%%%%%%%%%%%%%%%%%%%%%%%%%%%%%%%%%%%%%%%%%%%%%%%%%%%%%%%%%%

There are

\[
\binom nk
\]

ways to choose which \(k\) trials are successes.

Therefore,

\[
\boxed{
P(\text{exactly }k\text{ successes})
=
\binom nk
p^k(1-p)^{n-k}.
}
\]

This formula leads directly to the binomial distribution.

%%%%%%%%%%%%%%%%%%%%%%%%%%%%%%%%%%%%%%%%%%%%%%%%%%%%%%%%%%%%
\subsection{Worked Example: Exactly Three Heads}
\label{subsec:worked-three-heads}
%%%%%%%%%%%%%%%%%%%%%%%%%%%%%%%%%%%%%%%%%%%%%%%%%%%%%%%%%%%%

\begin{workedexamplebox}{Five Fair Coin Tosses}

Toss a fair coin five times.

Find the probability of exactly three heads.

There are

\[
\binom53
\]

ways to choose the positions of the heads.

Each particular sequence has probability

\[
\left(\frac12\right)^5.
\]

Therefore,

\[
P(X=3)
=
\binom53
\left(\frac12\right)^5.
\]

\begin{answerbox}
\[
\boxed{
P(X=3)
=
\frac{10}{32}
=
\frac5{16}.
}
\]
\end{answerbox}

\end{workedexamplebox}

%%%%%%%%%%%%%%%%%%%%%%%%%%%%%%%%%%%%%%%%%%%%%%%%%%%%%%%%%%%%
\subsection{Tree Diagrams}
\label{subsec:probability-tree-diagrams}
%%%%%%%%%%%%%%%%%%%%%%%%%%%%%%%%%%%%%%%%%%%%%%%%%%%%%%%%%%%%

A probability tree represents a sequence of conditional choices.

Each branch is labeled with a conditional probability.

To find the probability of one complete path, multiply along the path.

To combine mutually exclusive paths leading to the same event, add the
path probabilities.

Thus probability trees encode:

\[
\boxed{
\text{multiply along branches}
}
\]

and

\[
\boxed{
\text{add across disjoint paths}.
}
\]

%%%%%%%%%%%%%%%%%%%%%%%%%%%%%%%%%%%%%%%%%%%%%%%%%%%%%%%%%%%%
\subsection{Probability Tables}
\label{subsec:probability-tables}
%%%%%%%%%%%%%%%%%%%%%%%%%%%%%%%%%%%%%%%%%%%%%%%%%%%%%%%%%%%%

For two categorical events or classifications, a table can organize
joint probabilities.

For example,

\[
\begin{array}{c|cc|c}
 & B & B^c & \text{Total}\\
\hline
A & P(A\cap B) & P(A\cap B^c) & P(A)\\
A^c & P(A^c\cap B) & P(A^c\cap B^c) & P(A^c)\\
\hline
\text{Total} & P(B) & P(B^c) & 1
\end{array}
\]

Conditional probabilities can be obtained by dividing joint entries by
the appropriate row or column total.

%%%%%%%%%%%%%%%%%%%%%%%%%%%%%%%%%%%%%%%%%%%%%%%%%%%%%%%%%%%%
\subsection{Worked Example: Two-Way Probability Table}
\label{subsec:worked-probability-table}
%%%%%%%%%%%%%%%%%%%%%%%%%%%%%%%%%%%%%%%%%%%%%%%%%%%%%%%%%%%%

\begin{workedexamplebox}{Conditional Probability from a Table}

Suppose

\[
P(A\cap B)=0.24,
\]

\[
P(A\cap B^c)=0.16,
\]

\[
P(A^c\cap B)=0.36,
\]

and

\[
P(A^c\cap B^c)=0.24.
\]

Then

\[
P(B)=0.24+0.36=0.60.
\]

Therefore,

\[
P(A\mid B)
=
\frac{0.24}{0.60}.
\]

\begin{answerbox}
\[
\boxed{
P(A\mid B)=0.40.
}
\]
\end{answerbox}

\end{workedexamplebox}

%%%%%%%%%%%%%%%%%%%%%%%%%%%%%%%%%%%%%%%%%%%%%%%%%%%%%%%%%%%%
\subsection{Odds}
\label{subsec:probability-odds}
%%%%%%%%%%%%%%%%%%%%%%%%%%%%%%%%%%%%%%%%%%%%%%%%%%%%%%%%%%%%

Suppose an event has probability

\[
P(A)=p.
\]

The odds in favor of \(A\) are

\[
\boxed{
\frac{p}{1-p}.
}
\]

The odds against \(A\) are

\[
\boxed{
\frac{1-p}{p}.
}
\]

If the odds in favor are

\[
a:b,
\]

then

\[
\boxed{
P(A)
=
\frac{a}{a+b}.
}
\]

%%%%%%%%%%%%%%%%%%%%%%%%%%%%%%%%%%%%%%%%%%%%%%%%%%%%%%%%%%%%
\subsection{Worked Example: Odds to Probability}
\label{subsec:worked-odds}
%%%%%%%%%%%%%%%%%%%%%%%%%%%%%%%%%%%%%%%%%%%%%%%%%%%%%%%%%%%%

\begin{workedexamplebox}{Converting Odds}

Suppose the odds in favor of an event are

\[
3:2.
\]

Then

\[
P(A)
=
\frac{3}{3+2}.
\]

\begin{answerbox}
\[
\boxed{
P(A)=\frac35.
}
\]
\end{answerbox}

\end{workedexamplebox}

%%%%%%%%%%%%%%%%%%%%%%%%%%%%%%%%%%%%%%%%%%%%%%%%%%%%%%%%%%%%
\subsection{Classical Probability}
\label{subsec:classical-probability}
%%%%%%%%%%%%%%%%%%%%%%%%%%%%%%%%%%%%%%%%%%%%%%%%%%%%%%%%%%%%

The classical interpretation applies when outcomes are symmetric and
equally likely.

Then

\[
\boxed{
P(A)
=
\frac{
\text{number of favorable outcomes}
}{
\text{number of possible outcomes}
}.
}
\]

This works naturally for:

\begin{itemize}
    \item fair dice;
    \item shuffled cards;
    \item uniformly random permutations;
    \item uniformly selected subsets.
\end{itemize}

%%%%%%%%%%%%%%%%%%%%%%%%%%%%%%%%%%%%%%%%%%%%%%%%%%%%%%%%%%%%
\subsection{Frequentist Interpretation}
\label{subsec:frequentist-probability}
%%%%%%%%%%%%%%%%%%%%%%%%%%%%%%%%%%%%%%%%%%%%%%%%%%%%%%%%%%%%

The frequentist interpretation connects probability with long-run
relative frequency.

If an experiment is repeated many times under similar conditions, the
fraction of trials in which event \(A\) occurs may stabilize near

\[
P(A).
\]

Symbolically,

\[
\boxed{
\frac{
\text{number of occurrences of }A
}{
\text{number of trials}
}
\approx
P(A)
}
\]

for a large number of trials.

This intuition will later be formalized by laws of large numbers.

%%%%%%%%%%%%%%%%%%%%%%%%%%%%%%%%%%%%%%%%%%%%%%%%%%%%%%%%%%%%
\subsection{Subjective and Bayesian Probability}
\label{subsec:subjective-bayesian-probability}
%%%%%%%%%%%%%%%%%%%%%%%%%%%%%%%%%%%%%%%%%%%%%%%%%%%%%%%%%%%%

Probability may also represent quantified uncertainty about an unknown
state.

In Bayesian analysis, probabilities can represent uncertainty before
data are observed and are updated using Bayes' theorem.

Thus

\[
\boxed{
\text{prior information}
+
\text{data}
\longrightarrow
\text{posterior information}.
}
\]

The mathematical probability axioms remain the same regardless of the
interpretation.

%%%%%%%%%%%%%%%%%%%%%%%%%%%%%%%%%%%%%%%%%%%%%%%%%%%%%%%%%%%%
\subsection{Probability Models}
\label{subsec:probability-models}
%%%%%%%%%%%%%%%%%%%%%%%%%%%%%%%%%%%%%%%%%%%%%%%%%%%%%%%%%%%%

A complete probability model specifies:

\begin{enumerate}
    \item the sample space;
    \item the allowable events;
    \item a probability measure.
\end{enumerate}

In finite elementary problems, all subsets of \(\Omega\) may be events.

For uncountable spaces, a more careful collection of measurable events
is needed.

This motivates measure-theoretic probability later in the chapter.

%%%%%%%%%%%%%%%%%%%%%%%%%%%%%%%%%%%%%%%%%%%%%%%%%%%%%%%%%%%%
\subsection{Uniform Finite Probability}
\label{subsec:uniform-finite-probability}
%%%%%%%%%%%%%%%%%%%%%%%%%%%%%%%%%%%%%%%%%%%%%%%%%%%%%%%%%%%%

If \(\Omega\) is finite with

\[
|\Omega|=N
\]

and all outcomes are equally likely, then

\[
P(\{\omega\})
=
\frac1N
\]

for every

\[
\omega\in\Omega.
\]

Therefore,

\[
\boxed{
P(A)
=
\sum_{\omega\in A}P(\{\omega\})
=
\frac{|A|}{N}.
}
\]

%%%%%%%%%%%%%%%%%%%%%%%%%%%%%%%%%%%%%%%%%%%%%%%%%%%%%%%%%%%%
\subsection{Nonuniform Finite Probability}
\label{subsec:nonuniform-finite-probability}
%%%%%%%%%%%%%%%%%%%%%%%%%%%%%%%%%%%%%%%%%%%%%%%%%%%%%%%%%%%%

Finite outcomes need not be equally likely.

Suppose

\[
\Omega
=
\{\omega_1,\ldots,\omega_n\}
\]

with

\[
P(\{\omega_i\})=p_i.
\]

Then

\[
\boxed{
p_i\geq0
}
\]

and

\[
\boxed{
\sum_{i=1}^{n}p_i=1.
}
\]

For any event \(A\),

\[
\boxed{
P(A)
=
\sum_{\omega_i\in A}p_i.
}
\]

%%%%%%%%%%%%%%%%%%%%%%%%%%%%%%%%%%%%%%%%%%%%%%%%%%%%%%%%%%%%
\subsection{Worked Example: Loaded Die}
\label{subsec:worked-loaded-die}
%%%%%%%%%%%%%%%%%%%%%%%%%%%%%%%%%%%%%%%%%%%%%%%%%%%%%%%%%%%%

\begin{workedexamplebox}{Nonuniform Outcomes}

Suppose a loaded die has probabilities

\[
P(1)=P(2)=P(3)=P(4)=P(5)=0.15
\]

and

\[
P(6)=0.25.
\]

The probabilities sum to

\[
5(0.15)+0.25=1.
\]

The probability of rolling an even number is

\[
P(2)+P(4)+P(6).
\]

Thus,

\begin{answerbox}
\[
\boxed{
P(\text{even})
=
0.15+0.15+0.25
=
0.55.
}
\]
\end{answerbox}

\end{workedexamplebox}

%%%%%%%%%%%%%%%%%%%%%%%%%%%%%%%%%%%%%%%%%%%%%%%%%%%%%%%%%%%%
\subsection{At Least One Events}
\label{subsec:at-least-one-events}
%%%%%%%%%%%%%%%%%%%%%%%%%%%%%%%%%%%%%%%%%%%%%%%%%%%%%%%%%%%%

For repeated trials, probabilities involving

\[
\text{at least one}
\]

are often easier through complements.

If \(A_i\) is the event that success occurs on trial \(i\), then

\[
\left(
\bigcup_{i=1}^{n}A_i
\right)^c
=
\bigcap_{i=1}^{n}A_i^c.
\]

Therefore,

\[
\boxed{
P(\text{at least one success})
=
1-P(\text{no successes}).
}
\]

For independent identical trials,

\[
\boxed{
P(\text{at least one success})
=
1-(1-p)^n.
}
\]

%%%%%%%%%%%%%%%%%%%%%%%%%%%%%%%%%%%%%%%%%%%%%%%%%%%%%%%%%%%%
\subsection{At Most and At Least}
\label{subsec:at-most-at-least}
%%%%%%%%%%%%%%%%%%%%%%%%%%%%%%%%%%%%%%%%%%%%%%%%%%%%%%%%%%%%

Probability language must be interpreted carefully.

\[
\text{at most }k
\]

means

\[
\boxed{
X\leq k.
}
\]

\[
\text{at least }k
\]

means

\[
\boxed{
X\geq k.
}
\]

\[
\text{more than }k
\]

means

\[
\boxed{
X>k.
}
\]

\[
\text{fewer than }k
\]

means

\[
\boxed{
X<k.
}
\]

%%%%%%%%%%%%%%%%%%%%%%%%%%%%%%%%%%%%%%%%%%%%%%%%%%%%%%%%%%%%
\subsection{Exactly One of Two Events}
\label{subsec:exactly-one-event}
%%%%%%%%%%%%%%%%%%%%%%%%%%%%%%%%%%%%%%%%%%%%%%%%%%%%%%%%%%%%

The event that exactly one of \(A\) or \(B\) occurs is

\[
\boxed{
(A\cap B^c)
\cup
(A^c\cap B).
}
\]

The two pieces are disjoint.

Therefore,

\[
\boxed{
P(\text{exactly one})
=
P(A)+P(B)-2P(A\cap B).
}
\]

%%%%%%%%%%%%%%%%%%%%%%%%%%%%%%%%%%%%%%%%%%%%%%%%%%%%%%%%%%%%
\subsection{Neither Event}
\label{subsec:neither-event}
%%%%%%%%%%%%%%%%%%%%%%%%%%%%%%%%%%%%%%%%%%%%%%%%%%%%%%%%%%%%

The event that neither \(A\) nor \(B\) occurs is

\[
A^c\cap B^c.
\]

By De Morgan's law,

\[
A^c\cap B^c
=
(A\cup B)^c.
\]

Therefore,

\[
\boxed{
P(\text{neither})
=
1-P(A\cup B).
}
\]

%%%%%%%%%%%%%%%%%%%%%%%%%%%%%%%%%%%%%%%%%%%%%%%%%%%%%%%%%%%%
\subsection{Conditional Complements}
\label{subsec:conditional-complements}
%%%%%%%%%%%%%%%%%%%%%%%%%%%%%%%%%%%%%%%%%%%%%%%%%%%%%%%%%%%%

Given \(B\),

\[
A
\]

and

\[
A^c
\]

still form complementary events inside the conditional sample space.

Therefore,

\[
\boxed{
P(A^c\mid B)
=
1-P(A\mid B).
}
\]

This identity is often useful in sequential problems.

%%%%%%%%%%%%%%%%%%%%%%%%%%%%%%%%%%%%%%%%%%%%%%%%%%%%%%%%%%%%
\subsection{Conditional Addition Rule}
\label{subsec:conditional-addition-rule}
%%%%%%%%%%%%%%%%%%%%%%%%%%%%%%%%%%%%%%%%%%%%%%%%%%%%%%%%%%%%

Conditioning on \(C\) produces a new probability measure whenever

\[
P(C)>0.
\]

Thus,

\[
\boxed{
P(A\cup B\mid C)
=
P(A\mid C)
+
P(B\mid C)
-
P(A\cap B\mid C).
}
\]

Ordinary probability identities remain valid inside the conditional
sample space.

%%%%%%%%%%%%%%%%%%%%%%%%%%%%%%%%%%%%%%%%%%%%%%%%%%%%%%%%%%%%
\subsection{Independence of Complements}
\label{subsec:independence-complements}
%%%%%%%%%%%%%%%%%%%%%%%%%%%%%%%%%%%%%%%%%%%%%%%%%%%%%%%%%%%%

If \(A\) and \(B\) are independent, then so are:

\[
A^c\text{ and }B,
\]

\[
A\text{ and }B^c,
\]

and

\[
A^c\text{ and }B^c.
\]

For example,

\[
P(A^c\cap B)
=
P(B)-P(A\cap B).
\]

Using independence,

\[
P(A^c\cap B)
=
P(B)-P(A)P(B).
\]

Therefore,

\[
P(A^c\cap B)
=
(1-P(A))P(B)
=
P(A^c)P(B).
\]

%%%%%%%%%%%%%%%%%%%%%%%%%%%%%%%%%%%%%%%%%%%%%%%%%%%%%%%%%%%%
\subsection{Bernoulli Trials Preview}
\label{subsec:bernoulli-trials-preview}
%%%%%%%%%%%%%%%%%%%%%%%%%%%%%%%%%%%%%%%%%%%%%%%%%%%%%%%%%%%%

A Bernoulli trial has exactly two outcomes:

\[
\boxed{
\text{success}
}
\]

and

\[
\boxed{
\text{failure}.
}
\]

If

\[
P(\text{success})=p,
\]

then

\[
P(\text{failure})=1-p.
\]

Independent repetitions of a Bernoulli trial form the foundation of the
binomial and geometric distributions.

%%%%%%%%%%%%%%%%%%%%%%%%%%%%%%%%%%%%%%%%%%%%%%%%%%%%%%%%%%%%
\subsection{Probability of a Sequence of Bernoulli Outcomes}
\label{subsec:bernoulli-sequence}
%%%%%%%%%%%%%%%%%%%%%%%%%%%%%%%%%%%%%%%%%%%%%%%%%%%%%%%%%%%%

Suppose independent trials have success probability \(p\).

For a specific sequence such as

\[
SFSFSS,
\]

count the number of successes and failures.

If there are \(k\) successes and \(n-k\) failures, then the probability
of that exact sequence is

\[
\boxed{
p^k(1-p)^{n-k}.
}
\]

The locations of the successes matter when describing one specific
sequence.

%%%%%%%%%%%%%%%%%%%%%%%%%%%%%%%%%%%%%%%%%%%%%%%%%%%%%%%%%%%%
\subsection{Counting All Success Arrangements}
\label{subsec:count-success-arrangements}
%%%%%%%%%%%%%%%%%%%%%%%%%%%%%%%%%%%%%%%%%%%%%%%%%%%%%%%%%%%%

If only the number of successes matters, there are

\[
\binom nk
\]

possible positions for the \(k\) successes.

Hence,

\[
\boxed{
\binom nkp^k(1-p)^{n-k}.
}
\]

This will later be recognized as a probability mass function.

%%%%%%%%%%%%%%%%%%%%%%%%%%%%%%%%%%%%%%%%%%%%%%%%%%%%%%%%%%%%
\subsection{Probability and Combinatorial Symmetry}
\label{subsec:probability-combinatorial-symmetry}
%%%%%%%%%%%%%%%%%%%%%%%%%%%%%%%%%%%%%%%%%%%%%%%%%%%%%%%%%%%%

When all objects are selected uniformly from a finite set, symmetry
often determines probabilities immediately.

For example, if one element is selected uniformly from \([n]\), then

\[
\boxed{
P(\text{selected element is }i)
=
\frac1n.
}
\]

For a uniformly random permutation, every element is equally likely to
appear in every position.

Symmetry is therefore one of the most important tools for elementary
probability.

%%%%%%%%%%%%%%%%%%%%%%%%%%%%%%%%%%%%%%%%%%%%%%%%%%%%%%%%%%%%
\subsection{Random Permutations}
\label{subsec:elementary-random-permutations}
%%%%%%%%%%%%%%%%%%%%%%%%%%%%%%%%%%%%%%%%%%%%%%%%%%%%%%%%%%%%

A uniformly random permutation of \(n\) objects chooses each of the

\[
n!
\]

permutations with probability

\[
\boxed{
\frac1{n!}.
}
\]

For a fixed object, the probability of occupying any specified position
is

\[
\boxed{
\frac1n.
}
\]

This model appears throughout probability and combinatorics.

%%%%%%%%%%%%%%%%%%%%%%%%%%%%%%%%%%%%%%%%%%%%%%%%%%%%%%%%%%%%
\subsection{Worked Example: Random Seating}
\label{subsec:worked-random-seating}
%%%%%%%%%%%%%%%%%%%%%%%%%%%%%%%%%%%%%%%%%%%%%%%%%%%%%%%%%%%%

\begin{workedexamplebox}{Two People Sitting Together}

Suppose \(8\) people are seated uniformly at random in a row.

Find the probability that two specified people, \(A\) and \(B\), sit
next to one another.

There are

\[
8!
\]

total arrangements.

Treat \(A\) and \(B\) as one combined block.

Then there are

\[
7!
\]

ways to arrange this block with the other \(6\) people.

Inside the block, \(A\) and \(B\) can appear in

\[
2
\]

orders.

Thus the favorable count is

\[
2(7!).
\]

Therefore,

\begin{answerbox}
\[
\boxed{
P(A\text{ next to }B)
=
\frac{2(7!)}{8!}
=
\frac14.
}
\]
\end{answerbox}

\end{workedexamplebox}

%%%%%%%%%%%%%%%%%%%%%%%%%%%%%%%%%%%%%%%%%%%%%%%%%%%%%%%%%%%%
\subsection{Random Subsets}
\label{subsec:elementary-random-subsets}
%%%%%%%%%%%%%%%%%%%%%%%%%%%%%%%%%%%%%%%%%%%%%%%%%%%%%%%%%%%%

Suppose a \(k\)-element subset is selected uniformly from an
\(n\)-element set.

There are

\[
\binom nk
\]

possible subsets.

For a fixed element \(x\), the number of selected subsets containing
\(x\) is

\[
\binom{n-1}{k-1}.
\]

Thus,

\[
P(x\text{ selected})
=
\frac{
\binom{n-1}{k-1}
}{
\binom nk
}.
\]

Simplifying,

\[
\boxed{
P(x\text{ selected})
=
\frac{k}{n}.
}
\]

%%%%%%%%%%%%%%%%%%%%%%%%%%%%%%%%%%%%%%%%%%%%%%%%%%%%%%%%%%%%
\subsection{Worked Example: Random Committee}
\label{subsec:worked-random-committee}
%%%%%%%%%%%%%%%%%%%%%%%%%%%%%%%%%%%%%%%%%%%%%%%%%%%%%%%%%%%%

\begin{workedexamplebox}{Probability a Particular Person Is Selected}

A committee of \(5\) is chosen uniformly from \(20\) people.

For one particular person \(A\),

\[
P(A\text{ selected})
=
\frac{
\binom{19}{4}
}{
\binom{20}{5}
}.
\]

Using the random-subset identity,

\begin{answerbox}
\[
\boxed{
P(A\text{ selected})
=
\frac5{20}
=
\frac14.
}
\]
\end{answerbox}

\end{workedexamplebox}

%%%%%%%%%%%%%%%%%%%%%%%%%%%%%%%%%%%%%%%%%%%%%%%%%%%%%%%%%%%%
\subsection{Sampling Bias and Model Assumptions}
\label{subsec:probability-model-assumptions}
%%%%%%%%%%%%%%%%%%%%%%%%%%%%%%%%%%%%%%%%%%%%%%%%%%%%%%%%%%%%

A probability calculation is only as meaningful as its model.

For example, using

\[
P(A)
=
\frac{|A|}{|\Omega|}
\]

requires equally likely outcomes.

If the physical process is biased, this assumption may be false.

Therefore one should always ask:

\begin{itemize}
    \item What is the sample space?
    \item Which outcomes are equally likely?
    \item Is independence justified?
    \item Is replacement occurring?
    \item Is order relevant?
\end{itemize}

%%%%%%%%%%%%%%%%%%%%%%%%%%%%%%%%%%%%%%%%%%%%%%%%%%%%%%%%%%%%
\subsection{Common Probability Workflow}
\label{subsec:elementary-probability-workflow}
%%%%%%%%%%%%%%%%%%%%%%%%%%%%%%%%%%%%%%%%%%%%%%%%%%%%%%%%%%%%

A useful general workflow is:

\begin{enumerate}
    \item identify the experiment;
    \item define the sample space;
    \item define the relevant events;
    \item determine whether outcomes are equally likely;
    \item decide whether counting is needed;
    \item identify dependence or independence;
    \item use complements if they simplify the event;
    \item apply conditional probability when information is given;
    \item check that the final answer lies between \(0\) and \(1\).
\end{enumerate}

%%%%%%%%%%%%%%%%%%%%%%%%%%%%%%%%%%%%%%%%%%%%%%%%%%%%%%%%%%%%
\subsection{Common Errors}
\label{subsec:elementary-probability-common-errors}
%%%%%%%%%%%%%%%%%%%%%%%%%%%%%%%%%%%%%%%%%%%%%%%%%%%%%%%%%%%%

\subsubsection{Assuming Outcomes Are Equally Likely}

The formula

\[
P(A)
=
\frac{|A|}{|\Omega|}
\]

works only when all outcomes in \(\Omega\) are equally likely.

%%%%%%%%%%%%%%%%%%%%%%%%%%%%%%%%%%%%%%%%%%%%%%%%%%%%%%%%%%%%
\subsubsection{Confusing an Event with an Outcome}

An outcome is one element of the sample space.

An event is a set of outcomes.

%%%%%%%%%%%%%%%%%%%%%%%%%%%%%%%%%%%%%%%%%%%%%%%%%%%%%%%%%%%%
\subsubsection{Adding Probabilities Without Correcting Overlap}

In general,

\[
P(A\cup B)
\neq
P(A)+P(B).
\]

The correct formula is

\[
P(A\cup B)
=
P(A)+P(B)-P(A\cap B).
\]

%%%%%%%%%%%%%%%%%%%%%%%%%%%%%%%%%%%%%%%%%%%%%%%%%%%%%%%%%%%%
\subsubsection{Confusing Mutually Exclusive and Independent}

Mutually exclusive means the events cannot occur together.

Independent means occurrence of one does not change the probability of
the other.

These are not equivalent.

%%%%%%%%%%%%%%%%%%%%%%%%%%%%%%%%%%%%%%%%%%%%%%%%%%%%%%%%%%%%
\subsubsection{Multiplying Probabilities Without Checking Dependence}

In general,

\[
P(A\cap B)
\neq
P(A)P(B).
\]

Instead,

\[
P(A\cap B)
=
P(A)P(B\mid A).
\]

The simpler product applies only under independence.

%%%%%%%%%%%%%%%%%%%%%%%%%%%%%%%%%%%%%%%%%%%%%%%%%%%%%%%%%%%%
\subsubsection{Reversing Conditional Probability}

In general,

\[
\boxed{
P(A\mid B)
\neq
P(B\mid A).
}
\]

Bayes' theorem is needed to reverse the condition.

%%%%%%%%%%%%%%%%%%%%%%%%%%%%%%%%%%%%%%%%%%%%%%%%%%%%%%%%%%%%
\subsubsection{Ignoring the Base Rate}

A highly accurate test may still have a modest posterior probability
when the underlying event is rare.

Bayes' theorem automatically incorporates the base probability.

%%%%%%%%%%%%%%%%%%%%%%%%%%%%%%%%%%%%%%%%%%%%%%%%%%%%%%%%%%%%
\subsubsection{Confusing With Replacement and Without Replacement}

Without replacement, probabilities usually change from one draw to the
next.

With replacement, the original probabilities are restored.

%%%%%%%%%%%%%%%%%%%%%%%%%%%%%%%%%%%%%%%%%%%%%%%%%%%%%%%%%%%%
\subsubsection{Mixing Ordered and Unordered Counts}

If the denominator counts ordered outcomes, the numerator must use the
same convention.

If the denominator counts unordered outcomes, the numerator must also
be unordered.

%%%%%%%%%%%%%%%%%%%%%%%%%%%%%%%%%%%%%%%%%%%%%%%%%%%%%%%%%%%%
\subsubsection{Forgetting the Complement Shortcut}

For events involving ``at least one,'' calculating

\[
1-P(\text{none})
\]

is often substantially easier than direct enumeration.

%%%%%%%%%%%%%%%%%%%%%%%%%%%%%%%%%%%%%%%%%%%%%%%%%%%%%%%%%%%%
\subsection{Applications}
\label{subsec:elementary-probability-applications}
%%%%%%%%%%%%%%%%%%%%%%%%%%%%%%%%%%%%%%%%%%%%%%%%%%%%%%%%%%%%

\begin{applicationbox}

\textbf{Statistics.}
Probability provides the mathematical foundation for sampling,
estimation, confidence intervals, and hypothesis testing.

\medskip

\textbf{Reliability.}
Component failures and system reliability are modeled using events and
conditional probabilities.

\medskip

\textbf{Medicine.}
Bayes' theorem is essential for interpreting diagnostic-test results.

\medskip

\textbf{Engineering.}
Probability models measurement uncertainty, manufacturing variation,
and system risk.

\medskip

\textbf{Finance.}
Uncertain returns, defaults, losses, and market movements are modeled
probabilistically.

\medskip

\textbf{Computer Science.}
Randomized algorithms, hashing, machine learning, and network models
depend on probability.

\medskip

\textbf{Operations Research.}
Queues, inventories, scheduling systems, and stochastic optimization
require probabilistic models.

\medskip

\textbf{Biology.}
Genetics, epidemiology, population models, and experimental data use
probabilistic reasoning.

\medskip

\textbf{Combinatorics.}
Uniform random permutations, random subsets, and random graphs combine
counting with probability.

\end{applicationbox}

%%%%%%%%%%%%%%%%%%%%%%%%%%%%%%%%%%%%%%%%%%%%%%%%%%%%%%%%%%%%
\subsection{Exercises}
\label{subsec:elementary-probability-exercises}
%%%%%%%%%%%%%%%%%%%%%%%%%%%%%%%%%%%%%%%%%%%%%%%%%%%%%%%%%%%%

\begin{exercise}[1]
Define a random experiment.
\end{exercise}

\begin{exercise}[1]
Define an outcome.
\end{exercise}

\begin{exercise}[1]
Define a sample space.
\end{exercise}

\begin{exercise}[1]
Define an event.
\end{exercise}

\begin{exercise}[1]
For one fair die roll, write the sample space.
\end{exercise}

\begin{exercise}[1]
For two fair coin tosses, write the sample space.
\end{exercise}

\begin{exercise}[1]
For a fair die, find the probability of rolling:
\begin{itemize}
    \item an even number;
    \item a number greater than \(4\);
    \item a prime number.
\end{itemize}
\end{exercise}

\begin{exercise}[1]
State the three probability axioms.
\end{exercise}

\begin{exercise}[1]
Prove that

\[
P(\varnothing)=0.
\]
\end{exercise}

\begin{exercise}[1]
State the complement rule.
\end{exercise}

\begin{exercise}[1]
Suppose

\[
P(A)=0.37.
\]

Find

\[
P(A^c).
\]
\end{exercise}

\begin{exercise}[1]
Define mutually exclusive events.
\end{exercise}

\begin{exercise}[1]
Define independent events.
\end{exercise}

\begin{exercise}[1]
Explain why mutually exclusive events of positive probability cannot be
independent.
\end{exercise}

\begin{exercise}[2]
For events \(A\) and \(B\), suppose

\[
P(A)=0.6,
\qquad
P(B)=0.5,
\qquad
P(A\cap B)=0.3.
\]

Find

\[
P(A\cup B).
\]
\end{exercise}

\begin{exercise}[2]
Using the same values, determine whether \(A\) and \(B\) are
independent.
\end{exercise}

\begin{exercise}[2]
Suppose

\[
P(A)=0.7
\]

and

\[
P(B\mid A)=0.4.
\]

Find

\[
P(A\cap B).
\]
\end{exercise}

\begin{exercise}[2]
Suppose

\[
P(A\cap B)=0.24
\]

and

\[
P(B)=0.6.
\]

Find

\[
P(A\mid B).
\]
\end{exercise}

\begin{exercise}[2]
A fair die is rolled five times. Find the probability of obtaining no
sixes.
\end{exercise}

\begin{exercise}[2]
Use the previous result to find the probability of obtaining at least
one six.
\end{exercise}

\begin{exercise}[2]
A card is drawn from a standard deck. Find the probability that it is a
queen or a spade.
\end{exercise}

\begin{exercise}[2]
Two cards are drawn without replacement. Find the probability that both
are aces.
\end{exercise}

\begin{exercise}[2]
Repeat the previous problem with replacement.
\end{exercise}

\begin{exercise}[2]
A committee of \(4\) people is selected from \(12\). Find the probability
that a specified person is selected.
\end{exercise}

\begin{exercise}[2]
A committee of \(4\) is selected from \(12\). Find the probability that
two specified people are both selected.
\end{exercise}

\begin{exercise}[2]
Five fair coins are tossed. Find the probability of exactly two heads.
\end{exercise}

\begin{exercise}[2]
Five fair coins are tossed. Find the probability of at least one head.
\end{exercise}

\begin{exercise}[2]
Convert odds of

\[
4:3
\]

in favor of an event into probability.
\end{exercise}

\begin{exercise}[2]
An event has probability

\[
0.8.
\]

Find the odds in favor.
\end{exercise}

\begin{exercise}[3]
Prove the addition rule

\[
P(A\cup B)
=
P(A)+P(B)-P(A\cap B).
\]
\end{exercise}

\begin{exercise}[3]
Prove monotonicity:

\[
A\subseteq B
\Longrightarrow
P(A)\leq P(B).
\]
\end{exercise}

\begin{exercise}[3]
Derive the three-event inclusion--exclusion formula.
\end{exercise}

\begin{exercise}[3]
Prove the multiplication rule from the definition of conditional
probability.
\end{exercise}

\begin{exercise}[3]
Derive the sequential multiplication rule for three events.
\end{exercise}

\begin{exercise}[3]
Prove the law of total probability.
\end{exercise}

\begin{exercise}[3]
Derive Bayes' theorem from conditional probability.
\end{exercise}

\begin{exercise}[3]
A company receives \(70\%\) of its parts from supplier \(A\) and
\(30\%\) from supplier \(B\). Supplier \(A\) has defect rate \(1\%\),
while supplier \(B\) has defect rate \(4\%\). Find the probability that
a randomly selected part is defective.
\end{exercise}

\begin{exercise}[3]
For the previous problem, find the probability that a defective part
came from supplier \(B\).
\end{exercise}

\begin{exercise}[3]
Construct an example of three events that are pairwise independent but
not mutually independent.
\end{exercise}

\begin{exercise}[3]
Prove that if \(A\) and \(B\) are independent, then \(A^c\) and \(B\)
are independent.
\end{exercise}

\begin{exercise}[3]
Prove that if \(A\) and \(B\) are independent, then \(A^c\) and \(B^c\)
are independent.
\end{exercise}

\begin{exercise}[3]
A five-card hand is chosen from a standard deck. Find the probability
that all five cards have the same suit.
\end{exercise}

\begin{exercise}[3]
A five-card hand is chosen from a standard deck. Find the probability
of exactly two aces.
\end{exercise}

\begin{exercise}[3]
A group contains \(8\) women and \(7\) men. A committee of \(5\) is
chosen uniformly. Find the probability that exactly \(3\) women are
selected.
\end{exercise}

\begin{exercise}[3]
A box contains \(10\) good components and \(4\) defective components.
Five are selected without replacement. Find the probability of exactly
one defective component.
\end{exercise}

\begin{exercise}[3]
Explain why the previous problem has a hypergeometric structure.
\end{exercise}

\begin{exercise}[3]
A rare condition occurs in \(0.5\%\) of a population. A test has
sensitivity \(98\%\) and specificity \(97\%\). Compute the probability
that a person actually has the condition given a positive test.
\end{exercise}

\begin{exercise}[4]
Study the distinction between pairwise independence and mutual
independence for four or more events.
\end{exercise}

\begin{exercise}[4]
Investigate Simpson's paradox using conditional probability tables.
\end{exercise}

\begin{exercise}[4]
Study the prosecutor's fallacy and explain its relationship to reversing
conditional probabilities.
\end{exercise}

\begin{exercise}[4]
Investigate base-rate effects in diagnostic testing.
\end{exercise}

\begin{exercise}[4]
Study conditional independence and explain why

\[
A\perp B
\]

does not necessarily imply

\[
A\perp B\mid C.
\]
\end{exercise}

\begin{exercise}[4]
Investigate the Monty Hall problem using conditional probability.
\end{exercise}

\begin{exercise}[4]
Study the birthday problem and determine the probability that at least
two people share a birthday in a group of \(n\) people.
\end{exercise}

\begin{exercise}[4]
Investigate the coupon collector problem as a preview of discrete random
variables.
\end{exercise}

\begin{exercise}[4]
Study Bertrand's paradox and explain why continuous probability requires
care in specifying the probability model.
\end{exercise}

\begin{exercise}[4]
Investigate how finite probability spaces generalize to probability
measures on sigma-algebras.
\end{exercise}

%%%%%%%%%%%%%%%%%%%%%%%%%%%%%%%%%%%%%%%%%%%%%%%%%%%%%%%%%%%%
\subsection{Conceptual Summary}
\label{subsec:elementary-probability-summary}
%%%%%%%%%%%%%%%%%%%%%%%%%%%%%%%%%%%%%%%%%%%%%%%%%%%%%%%%%%%%

Probability begins with a sample space

\[
\boxed{
\Omega
}
\]

containing all possible outcomes.

An event is a subset

\[
\boxed{
A\subseteq\Omega.
}
\]

A probability measure satisfies

\[
\boxed{
P(A)\geq0,
}
\]

\[
\boxed{
P(\Omega)=1,
}
\]

and countable additivity over disjoint events.

For every event,

\[
\boxed{
0\leq P(A)\leq1.
}
\]

The complement rule is

\[
\boxed{
P(A^c)=1-P(A).
}
\]

For two events,

\[
\boxed{
P(A\cup B)
=
P(A)+P(B)-P(A\cap B).
}
\]

Conditional probability is

\[
\boxed{
P(A\mid B)
=
\frac{P(A\cap B)}{P(B)}.
}
\]

The multiplication rule is

\[
\boxed{
P(A\cap B)
=
P(A\mid B)P(B).
}
\]

Independence means

\[
\boxed{
P(A\cap B)=P(A)P(B).
}
\]

The law of total probability is

\[
\boxed{
P(A)
=
\sum_i
P(A\mid B_i)P(B_i).
}
\]

Bayes' theorem is

\[
\boxed{
P(B_j\mid A)
=
\frac{
P(A\mid B_j)P(B_j)
}{
\sum_iP(A\mid B_i)P(B_i)
}.
}
\]

For finite equally likely outcomes,

\[
\boxed{
P(A)
=
\frac{|A|}{|\Omega|}.
}
\]

For repeated independent Bernoulli trials,

\[
\boxed{
P(\text{exactly }k\text{ successes})
=
\binom nk
p^k(1-p)^{n-k}.
}
\]

Elementary probability therefore combines

\[
\boxed{
\text{sets}
+
\text{counting}
+
\text{conditional reasoning}
+
\text{uncertainty}.
}
\]

%%%%%%%%%%%%%%%%%%%%%%%%%%%%%%%%%%%%%%%%%%%%%%%%%%%%%%%%%%%%
\subsection{Section Connections}
\label{subsec:elementary-probability-connections}
%%%%%%%%%%%%%%%%%%%%%%%%%%%%%%%%%%%%%%%%%%%%%%%%%%%%%%%%%%%%

\chapterconnection{
Elementary Probability establishes the axiomatic and computational
foundation for Chapter 7. The next section, Random Variables, converts
numerical outcomes of random experiments into functions on the sample
space. Discrete and Continuous Distributions then develop standard
probability models. Joint Distributions extend conditional probability
and independence to multiple random variables. Expectation and Moments
generalize weighted averages, while Probability Transforms encode
distributions algebraically. The Laws of Large Numbers and Central Limit
Theorems formalize long-run frequency and approximation phenomena.
Measure-Theoretic Probability later places these ideas in their full
mathematical framework. Chapter 8 Statistics uses probability as the
foundation for inference from data.
}

%%%%%%%%%%%%%%%%%%%%%%%%%%%%%%%%%%%%%%%%%%%%%%%%%%%%%%%%%%%%
% SECTION SOURCE TRACKING
%%%%%%%%%%%%%%%%%%%%%%%%%%%%%%%%%%%%%%%%%%%%%%%%%%%%%%%%%%%%

%------------------------------------------------------------
% 7.1 Elementary Probability
%------------------------------------------------------------
%
% Source basis:
%   - Standard elementary probability theory.
%   - Standard axiomatic probability.
%   - Standard finite combinatorial probability.
%
% Used for:
%   - random experiments;
%   - outcomes;
%   - sample spaces;
%   - events;
%   - event set operations;
%   - Kolmogorov axioms;
%   - finite uniform probability;
%   - nonuniform finite probability;
%   - complements;
%   - monotonicity;
%   - addition rule;
%   - inclusion--exclusion;
%   - mutually exclusive events;
%   - conditional probability;
%   - multiplication rule;
%   - sequential probability;
%   - sampling with and without replacement;
%   - independence;
%   - pairwise and mutual independence;
%   - partitions of sample spaces;
%   - law of total probability;
%   - Bayes' theorem;
%   - priors and posteriors;
%   - diagnostic testing;
%   - sensitivity and specificity;
%   - combinatorial probability;
%   - ordered and unordered sampling;
%   - hypergeometric probability preview;
%   - repeated Bernoulli trials;
%   - binomial probability preview;
%   - probability trees;
%   - probability tables;
%   - odds;
%   - classical, frequentist, and Bayesian interpretations;
%   - at-least-one probabilities;
%   - random permutations and subsets;
%   - probability-model assumptions;
%   - applications and exercises.
%
% Notes:
%   - Random variables are introduced formally in Section 7.2.
%   - Binomial and hypergeometric distributions are developed in the
%     discrete-distribution section.
%   - General sigma-algebras and probability measures are developed
%     later in Measure-Theoretic Probability.
%   - Statistical applications are developed systematically in
%     Chapter 8.
%
%%%%%%%%%%%%%%%%%%%%%%%%%%%%%%%%%%%%%%%%%%%%%%%%%%%%%%%%%%%%